As the dial is vertical, the rod lies on the east-zenith-west plane. The
clock works only when the Sun is in the northern hemisphere,
$\delta_{\mathrm{Sun}} > \delta_{\mathrm{zenith}}$, because only in this case
there will be a shadow projected over the dial.

% FIGURE NEEDED: celestial sphere with zenith Z, horizon, SCP and NCP, local meridian, first vertical, and the Sun's diurnal circles at the summer solstice (south), equinox and winter solstice (south) (2012 theory general solutions, p. S1)

\begin{description}
\item[(i)] When $\delta_{\mathrm{Sol}} \le 0^\circ$, the Sun does not project
  a shadow when it is near the horizon (after rising and before setting). It
  corresponds to the time between the equinoxes: between September and March,
  during spring and summer.
  \hfill(3 pt declination; 2 pt months \& seasons)
\item[(ii)] When $\delta_{\mathrm{Sun}} \le -22^\circ 54'$, the solar clock
  doesn't work at all. It corresponds to the small time interval close to the
  Summer solstice, during December and January.
  \hfill(3 pt declination; 2 pt months \& seasons)
\end{description}
